\chapter{子列与 Bolzano--Weierstrass 定理：习题解答}

\begin{chapterguide}
本章解答特别区分“无穷多个指标”和“无穷多个不同取值”。抽取步骤均保证新指标严格递增。
\end{chapterguide}

\section*{A组　子列定义与继承}
\addcontentsline{toc}{section}{A组　子列定义与继承}

\begin{solution}{11.1}
对 \(k\) 归纳。\(n_1\ge1\)。若 \(n_k\ge k\)，严格递增性和整数性给出
\(n_{k+1}\ge n_k+1\ge k+1\)。故 \(n_k\ge k\)。给定 \(M\)，取整数 \(K>M\)，则 \(k\ge K\) 时 \(n_k\ge k>M\)，所以 \(n_k\to+\infty\)。
\end{solution}

\begin{solution}{11.2}
\(2k\) 与 \(k^2\) 都是严格递增正整数指标，故给出子列。\(2-k\) 很快不再是正整数，不能作为全体 \(k\) 的指标。\(\lfloor\sqrt k\rfloor+1\) 在许多相邻 \(k\) 上重复，不严格递增，所以一般不构成子列。
\end{solution}

\begin{solution}{11.3}
若 \(n_k\) 严格递增，\(k_j\) 也严格递增，则
\[
 j_1<j_2\Longrightarrow k_{j_1}<k_{j_2}
 \Longrightarrow n_{k_{j_1}}<n_{k_{j_2}}.
\]
复合指标严格递增。
\end{solution}

\begin{solution}{11.4}
有限极限情形按正文：原列从 \(N\) 起进入误差邻域，而 \(n_k\ge k\)，所以子列从第 \(N\) 项起也进入。对 \(+\infty\)，把误差邻域换成阈值 \(M\)；对 \(-\infty\) 对称。
\end{solution}

\begin{solution}{11.5}
反设原列不收敛到 \(L\)。由严格否定可抽出子列 \((a_{n_k})\)，使
\(\abs{a_{n_k}-L}\ge\varepsilon_0>0\) 恒成立。按假设该子列有进一步子列收敛到 \(L\)，但极限定义要求其误差最终小于 \(\varepsilon_0/2\)，矛盾。
\end{solution}

\begin{solution}{11.6}
不收敛到 \(L\) 给出
\[
 \exists\varepsilon_0>0\ \forall N\ \exists n\ge N,\quad
 \abs{a_n-L}\ge\varepsilon_0.
\]
先以 \(N=1\) 选 \(n_1\)。选好 \(n_k\) 后，以 \(N=n_k+1\) 选 \(n_{k+1}\)。所得指标严格递增，且每项都满足所需下界。反向由任一尾部均含该子列的后续项立即得到。
\end{solution}

\section*{B组　抽取与紧致性原型}
\addcontentsline{toc}{section}{B组　抽取与紧致性原型}

\begin{solution}{11.7}
偶数项为 \(1+1/(2k)\to1\)，奇数项为
\(-1+1/(2k-1)\to-1\)。两条子列极限不同，故原列发散。
\end{solution}

\begin{solution}{11.8}
若两个半区间都只含有限多个指标，则它们的指标集合之并有限；但每个原数列项都在至少一个半区间中，故全部自然数指标有限，矛盾。选 \(n_k\) 时，\(I_k\) 含无穷多个指标，删去有限集合 \(\{1,\ldots,n_{k-1}\}\) 后仍非空，所以可选更大的指标。
\end{solution}

\begin{solution}{11.9}
按二分区间构造 \(I_k\) 并选 \(a_{n_k}\in I_k\)。对 \(j,\ell\ge k\)，嵌套性给出两项都在 \(I_k\)，所以
\[
 \abs{a_{n_j}-a_{n_\ell}}
 \le\operatorname{length}(I_k)=2^{-k}(B-A).
\]
右端趋零，故子列为 Cauchy 列。
\end{solution}

\begin{solution}{11.10}
从无限有界集合 \(E\) 中递归选互异点 \(x_n\)。所得数列有收敛子列
\(x_{n_k}\to x\)。任意 \(x\) 的邻域含该子列的无穷多个后项；因所选点互异，其中可取不同于 \(x\) 的 \(E\) 中点，所以 \(x\) 是 \(E\) 的聚点。
\end{solution}

\begin{solution}{11.11}
若峰指标 \(n_1<n_2<\cdots\) 无穷，则因 \(n_{k+1}>n_k\)，峰性质给出
\(a_{n_k}\ge a_{n_{k+1}}\)。若峰指标有限，取 \(N\) 越过它们。每个 \(n\ge N\) 都不是峰指标，故有 \(m>n\) 满足 \(a_m>a_n\)。递归选择即得严格递增子列。
\end{solution}

\begin{solution}{11.12}
有界数列由单调子列定理得到一个单调子列；子列仍有界。单调有界收敛定理说明该子列收敛，故 Bolzano--Weierstrass 定理成立。
\end{solution}

\begin{solution}{11.13}
设唯一部分极限为 \(L\)。若原列不收敛到 \(L\)，可抽取始终远离 \(L\) 的子列。该子列仍有界，由 Bolzano--Weierstrass 定理有收敛的进一步子列，设极限为 \(M\)。距离下界传给极限，得 \(\abs{M-L}\ge\varepsilon_0\)，出现第二个部分极限，矛盾。
\end{solution}

\begin{solution}{11.14}
取
\[
 a_n=\begin{cases}0,&n\text{ 为偶数},\\ n,&n\text{ 为奇数}.\end{cases}
\]
它无界，但偶数子列恒为 \(0\) 并收敛。Bolzano--Weierstrass 只断言“有界蕴含存在收敛子列”，没有断言存在收敛子列蕴含有界。
\end{solution}

\begin{solution}{11.15}
有限多个值对应有限个指标集合，它们的并是 \(\N\)。若每个指标集合都有限，其并也有限，矛盾；故某个值出现无穷多次，产生常值子列。原列收敛当且仅当除某一个值外其余值都只出现有限次；此时原列最终恒于该值。必要性可取小于任意两个不同值距离一半的误差。
\end{solution}

\section*{C组　结构与项目}
\addcontentsline{toc}{section}{C组　结构与项目}

\begin{solution}{11.16}
定义 \(a_{2k}=a\)、\(a_{2k-1}=b\)。常值奇偶子列说明 \(a,b\) 都是部分极限。任意收敛子列若含无穷多个奇项或偶项，进一步子列恒为相应值；极限唯一性迫使其极限为 \(a\) 或 \(b\)。故部分极限恰为这两点。
\end{solution}

\begin{solution}{11.17}
枚举 \(\Q\cap[0,1]=\{q_1,q_2,\ldots\}\)，依次连接有限块
\[
 q_1;\quad q_1,q_2;\quad\ldots;\quad q_1,\ldots,q_m;\quad\cdots.
\]
对任意 \(x\in[0,1]\)，可选有理数 \(q_{j_k}\to x\)，并在足够靠后的不同块中各取一次，得到趋于 \(x\) 的子列。反之，所有项都在闭区间 \([0,1]\)，任意子列极限仍在该区间。因此部分极限集合恰为 \([0,1]\)。
\end{solution}

\begin{solution}{11.18}
从有界区间 \(I_0=[A,B]\) 开始，每步选含无穷多项的闭半区间 \(I_k\)。区间套给出交点 \(x\)。递归选 \(n_k>n_{k-1}\) 且 \(a_{n_k}\in I_k\)，于是
\[
 \abs{a_{n_k}-x}\le\operatorname{length}(I_k)
 =2^{-k}(B-A).
\]
该显式界直接证明子列收敛，并给出达到误差 \(\varepsilon\) 时只需
\(2^{-k}(B-A)<\varepsilon\)。
\end{solution}

\begin{solution}{11.19}
实数全序允许“峰指标”二分法，故总能抽出递增或递减子列。一般偏序中可能有无限反链：任意两项都不可比较，此时不存在长度大于一的单调子列。例如把互不包含的单点集按包含关系排序即可。
\end{solution}

\begin{solution}{11.20}
用三角形块构造：第 \(k\) 块以步长 \(1/k\) 从 \(-1\) 上升到 \(1\)，再下降到 \(-1\)。相邻差趋零，但每块都含 \(-1,1\)，故有两条分别收敛到这两个值的子列。答案是否定的。
\end{solution}

\begin{solution}{11.21}
正向：若原列趋于 \(L\)，任何子列本身已趋于 \(L\)。反向：若原列不趋于 \(L\)，抽出始终远离 \(L\) 的子列；该子列按假设还应含趋于 \(L\) 的进一步子列，与固定距离下界矛盾。
\end{solution}

\begin{solution}{11.22}
若值集 \(V=\{a_n\}\) 有聚点 \(x\)，每个 \(1/k\) 邻域含 \(V\) 中不同于 \(x\) 的点，并可递归选择出现指标严格增大的项，得到子列趋于 \(x\)。反向可能失败：常值列 \(a_n=c\) 有部分极限 \(c\)，但单点值集 \(\{c\}\) 在要求删心邻域含别点的定义下没有聚点。重复指标造成这一区别。
\end{solution}

\begin{solution}{11.23}
对有限集 \(F\)，把每个点周期性重复即可。对集合
\(F=\{0\}\cup\{1/m:m\in\N\}\)，可依次连接前 \(k\) 个非零点并插入 \(0\)，使每个点出现无穷次，同时所有非零点只在 \(0\) 附近累积。一般可数闭集若枚举为 \(\{x_j\}\)，连接有限块
\(x_1,\ldots,x_k\) 可保证每个枚举点为部分极限；闭性保证任何新出现的子列极限仍留在集合中。对无界集合需再安排有界窗口，避免某一阶段的远点妨碍局部抽取。
\end{solution}
